\documentclass[11pt]{article}
\usepackage[utf8]{inputenc}
\usepackage[T1]{fontenc}
\usepackage{url,hyperref,microtype}
\usepackage{booktabs}
\usepackage{amsmath}
\usepackage[onehalfspacing]{setspace}
\usepackage[margin=1in]{geometry}
\usepackage{natbib}

\title{Supplementary Materials\\[4pt]
\large Reading skill and the time course of morphological processing: electrophysiological evidence for stage-specific modulation}
\author{}
\date{}

\begin{document}
\maketitle
\thispagestyle{empty}

\section*{Overview}

This document reports the full model output for the word morphological family size analyses referenced in the main text. In the main manuscript, base frequency is treated as the primary morphological predictor for word ERP analyses, since the double dissociation central to the study's aims --- Orthographic Sensitivity (OS) modulating the N250 and Language Proficiency (LP) modulating the N400 --- is formally tested using base frequency as the morphological variable of interest. Family size analyses for words were conducted in parallel, using the same model structure, and are summarised briefly in the main text (see \emph{Words N250: Base Frequency} and \emph{Words N400: Base Frequency}). The full statistical output for these two family size models is reported below for transparency and to allow readers to evaluate the pattern in full.

\section*{S1. Words N250: Family Size}

Mean N250 amplitudes for word stimuli were analyzed using a linear mixed-effects model with fixed effects of Family Size, Orthographic Sensitivity (OS), Language Proficiency (LP), and the interaction between Family Size and OS. The model included random intercepts for participants and for electrode site nested within participant. The analysis revealed a robust main effect of Family Size---words from smaller morphological families elicited more negative N250 amplitudes than words from larger families, $M_{\mathrm{large}} = -0.57\,\mu$V, 95\% CI $[-1.03, -0.11]$, $M_{\mathrm{small}} = -0.90\,\mu$V, 95\% CI $[-1.36, -0.44]$, $\beta = 0.324$, $SE = 0.058$, $t(1618) = 5.54$, $p < .001$, $\eta^2_p = .02$. Neither Orthographic Sensitivity, $\beta = 0.261$, $SE = 0.256$, $t(57) = 1.02$, $p = .31$, nor Language Proficiency, $\beta = -0.279$, $SE = 0.187$, $t(57) = -1.50$, $p = .14$, yielded reliable main effects. The interaction between Family Size and Orthographic Sensitivity did not reach significance but showed a marginal trend, $\beta = -0.111$, $SE = 0.065$, $t(1618) = -1.70$, $p = .089$, $\eta^2_p = .002$, suggesting that the family size effect at the N250 may be modulated by orthographic skill, though this requires cautious interpretation given the marginal status of the effect. The model explained 1.9\% of the variance in N250 amplitude via fixed effects ($R^2_{\mathrm{marginal}} = .019$) and 73.8\% overall ($R^2_{\mathrm{conditional}} = .738$), with the large conditional $R^2$ reflecting substantial variability in N250 amplitude across participants and electrode sites. Together, these results indicate that morphological family structure robustly modulates early morpho-orthographic processing for words, with limited but suggestive evidence that orthographic skill may shape the sensitivity of this effect.

\begin{table}[h!]
	\centering
	\small
	\caption{Words N250: Family Size model, fixed effects.}
	\label{tab:s1}
	\begin{tabular}{lccccc}
		\toprule
		\textbf{Term} & $\beta$ & \textbf{SE} & \textbf{df} & \textbf{t} & \textbf{p} \\
		\midrule
		Family Size & 0.324 & 0.058 & 1618 & 5.54 & $< .001$ \\
		Orthographic Sensitivity & 0.261 & 0.256 & 57 & 1.02 & .31 \\
		Language Proficiency & $-0.279$ & 0.187 & 57 & $-1.50$ & .14 \\
		Family Size $\times$ OS & $-0.111$ & 0.065 & 1618 & $-1.70$ & .089 \\
		\bottomrule
	\end{tabular}
	\\[4pt]
	\raggedright \footnotesize $R^2_{\mathrm{marginal}} = .019$; $R^2_{\mathrm{conditional}} = .738$. Random effects: participant, electrode nested within participant.
\end{table}

\section*{S2. Words N400: Family Size}

Mean N400 amplitudes for word stimuli were analyzed using a linear mixed-effects model with fixed effects of Family Size, Language Proficiency (LP), Orthographic Sensitivity (OS), and the interaction between Family Size and LP. The model included random intercepts for participants and for electrode site nested within participant. The analysis revealed a robust main effect of Family Size---words from smaller morphological families elicited greater N400 amplitudes than words from larger families, $M_{\mathrm{large}} = 0.518\,\mu$V, 95\% CI $[-0.054, 1.089]$, $M_{\mathrm{small}} = 0.062\,\mu$V, 95\% CI $[-0.509, 0.634]$, $\beta = 0.452$, $SE = 0.064$, $t(1618) = 7.07$, $p < .001$, $\eta^2_p = .03$, indicating facilitated semantic integration for words supported by richer morphological families. A significant main effect of Orthographic Sensitivity was also observed, $\beta = 0.677$, $SE = 0.318$, $t(57) = 2.13$, $p = .038$, $\eta^2_p = .07$, reflecting overall differences in N400 amplitude as a function of orthographic skill independently of morphological structure. Language Proficiency did not yield a reliable main effect, $\beta = 0.001$, $SE = 0.231$, $t(57) = 0.005$, $p = .996$, nor did the Family Size $\times$ LP interaction reach significance, $\beta = -0.071$, $SE = 0.052$, $t(1618) = -1.37$, $p = .172$. These results indicate that semantic integration of morphological family information for familiar words is robust and not selectively modulated by individual differences in language proficiency.

\begin{table}[h!]
	\centering
	\small
	\caption{Words N400: Family Size model, fixed effects.}
	\label{tab:s2}
	\begin{tabular}{lccccc}
		\toprule
		\textbf{Term} & $\beta$ & \textbf{SE} & \textbf{df} & \textbf{t} & \textbf{p} \\
		\midrule
		Family Size & 0.452 & 0.064 & 1618 & 7.07 & $< .001$ \\
		Language Proficiency & 0.001 & 0.231 & 57 & 0.005 & .996 \\
		Orthographic Sensitivity & 0.677 & 0.318 & 57 & 2.13 & .038 \\
		Family Size $\times$ LP & $-0.071$ & 0.052 & 1618 & $-1.37$ & .172 \\
		\bottomrule
	\end{tabular}
\end{table}

\bibliographystyle{Frontiers-Harvard}
\bibliography{M21_LDT.bib}

\end{document}
